\begin{frame}{أمثلة تطبيقية}
    \begin{tbox}{\faLaptopCode\quad عشرة أمثلة من مجالات مختلفة}
        نطبق ما تعلمناه في مجالات من الحياة اليومية:
        الصحة، الرياضة، الطقس، النقل، التجارة الإلكترونية، الهواتف الذكية، الشبكات الاجتماعية، الطاقة، البيئة، البث الرقمي.
    \end{tbox}
\end{frame}

\begin{frame}[t]{المثال الأول: الصحة}
    \begin{columns}[T]
        \column{0.64\textwidth}
        \ccode[basicstyle=\codefont\codesmall]{d1/ex01.c}
        \column{0.33\textwidth}
        {\small\textbf{مؤشر كتلة الجسم}: نحسب مؤشر كتلة الجسم: الوزن مقسوم على مربع الطول.\par}
        \vspace{0.3cm}
        \out[basicstyle=\codefont\scriptsize\color{brandDark}]{d1/ex01.out}
    \end{columns}
\end{frame}

\begin{frame}[t]{المثال الثاني: الرياضة}
    \begin{columns}[T]
        \column{0.64\textwidth}
        \ccode[basicstyle=\codefont\codesmall]{d1/ex02.c}
        \column{0.33\textwidth}
        {\small\textbf{سرعة عدّاء}: عدّاء يقطع 10 كم في 52 دقيقة: نحسب سرعته ووتيرته.\par}
        \vspace{0.3cm}
        \out[basicstyle=\codefont\scriptsize\color{brandDark}]{d1/ex02.out}
    \end{columns}
\end{frame}

\begin{frame}[t]{المثال الثالث: الطقس}
    \begin{columns}[T]
        \column{0.64\textwidth}
        \ccode[basicstyle=\codefont\codesmall]{d1/ex03.c}
        \column{0.33\textwidth}
        {\small\textbf{تحويل درجة الحرارة}: نحوّل درجة الحرارة من سلسيوس إلى فهرنهايت.\par}
        \vspace{0.3cm}
        \out[basicstyle=\codefont\scriptsize\color{brandDark}]{d1/ex03.out}
    \end{columns}
\end{frame}

\begin{frame}[t]{المثال الرابع: النقل}
    \begin{columns}[T]
        \column{0.64\textwidth}
        \ccode[basicstyle=\codefont\codesmall]{d1/ex04.c}
        \column{0.33\textwidth}
        {\small\textbf{استهلاك الوقود}: نحسب استهلاك سيارة في 100 كم وتكلفة الرحلة.\par}
        \vspace{0.3cm}
        \out[basicstyle=\codefont\scriptsize\color{brandDark}]{d1/ex04.out}
    \end{columns}
\end{frame}

\begin{frame}[t]{المثال الخامس: التجارة الإلكترونية}
    \begin{columns}[T]
        \column{0.64\textwidth}
        \ccode[basicstyle=\codefont\codesmall]{d1/ex05.c}
        \column{0.33\textwidth}
        {\small\textbf{السعر بعد التخفيض}: نطبق تخفيضًا بنسبة مئوية ثم نضيف ثمن التوصيل.\par}
        \vspace{0.3cm}
        \out[basicstyle=\codefont\scriptsize\color{brandDark}]{d1/ex05.out}
    \end{columns}
\end{frame}

\begin{frame}[t]{المثال السادس: الهواتف الذكية}
    \begin{columns}[T]
        \column{0.64\textwidth}
        \ccode[basicstyle=\codefont\codesmall]{d1/ex06.c}
        \column{0.33\textwidth}
        {\small\textbf{مدة البطارية}: نحسب عدد ساعات عمل الهاتف من سعة البطارية والاستهلاك.\par}
        \vspace{0.3cm}
        \out[basicstyle=\codefont\scriptsize\color{brandDark}]{d1/ex06.out}
    \end{columns}
\end{frame}

\begin{frame}[t]{المثال السابع: الشبكات الاجتماعية}
    \begin{columns}[T]
        \column{0.64\textwidth}
        \ccode[basicstyle=\codefont\codesmall]{d1/ex07.c}
        \column{0.33\textwidth}
        {\small\textbf{معدل التفاعل}: نسبة التفاعل = (الإعجابات + التعليقات) مقسومة على المتابعين.\par}
        \vspace{0.3cm}
        \out[basicstyle=\codefont\scriptsize\color{brandDark}]{d1/ex07.out}
    \end{columns}
\end{frame}

\begin{frame}[t]{المثال الثامن: الطاقة}
    \begin{columns}[T]
        \column{0.64\textwidth}
        \ccode[basicstyle=\codefont\codesmall]{d1/ex08.c}
        \column{0.33\textwidth}
        {\small\textbf{فاتورة الكهرباء}: مدفأة استطاعتها 1500 واط تعمل 6.5 ساعات: نحسب الطاقة والتكلفة.\par}
        \vspace{0.3cm}
        \out[basicstyle=\codefont\scriptsize\color{brandDark}]{d1/ex08.out}
    \end{columns}
\end{frame}

\begin{frame}[t]{المثال التاسع: البيئة}
    \begin{columns}[T]
        \column{0.64\textwidth}
        \ccode[basicstyle=\codefont\codesmall]{d1/ex09.c}
        \column{0.33\textwidth}
        {\small\textbf{البصمة الكربونية}: نحسب انبعاثات سيارة في شهر، وعدد الأشجار اللازمة لامتصاصها في سنة.\par}
        \vspace{0.3cm}
        \out[basicstyle=\codefont\scriptsize\color{brandDark}]{d1/ex09.out}
    \end{columns}
\end{frame}

\begin{frame}[t]{المثال العاشر: البث الرقمي}
    \begin{columns}[T]
        \column{0.64\textwidth}
        \ccode[basicstyle=\codefont\codesmall]{d1/ex10.c}
        \column{0.33\textwidth}
        {\small\textbf{مدة فيديو}: نحوّل مدة فيديو بالثواني إلى ساعات ودقائق وثوانٍ، باستعمال \code{/} و\code{\%}.\par}
        \vspace{0.3cm}
        \out[basicstyle=\codefont\scriptsize\color{brandDark}]{d1/ex10.out}
    \end{columns}
\end{frame}
